\thispagestyle{first}
\begingroup
\setlength{\parindent}{0pt}
\begin{minipage}[c]{0.6\textwidth}
  \href{https://phaistoslabs.com}{\raisebox{-0.32\height}{\includegraphics[height=8.5mm]{figures/phaistos_logo.png}}%
  \hspace{2.6mm}{\heading\fontsize{13.5}{16}\selectfont\color{ink}Phaistos Labs}}
\end{minipage}%
\begin{minipage}[c]{0.4\textwidth}
  \raggedleft{\headingregular\itshape\small\color{ink2}Working paper\quad·\quad 11 October 2026}
\end{minipage}

\vspace{3.5mm}
{\color{rulecol}\rule{\textwidth}{0.6pt}}

\vspace{9mm}
{\headingregular\itshape\normalsize\color{ochretext}Preprint, not peer reviewed}

\vspace{2.5mm}
{\heading\fontsize{23}{27.5}\selectfont\color{ink}\raggedright
Liabut: a pre-registered computational test of how Timor’s Austronesian languages branch\par}

\vspace{5mm}
{\headingregular\large\color{ink}Dennis Crowley}\\[1mm]
{\headingregular\itshape\normalsize\color{ink2}Phaistos Labs}

\vspace{8mm}
\colorbox{paperlight}{\parbox{\dimexpr\textwidth-2\fboxsep\relax}{%
\vspace{2.5mm}
\hspace*{3mm}\parbox{\dimexpr\textwidth-2\fboxsep-6mm\relax}{\linespread{1.08}\selectfont\small
{\heading\color{ink}Abstract}\par\vspace{1.5mm}
Two classifications of Timor's Austronesian languages disagree on where Lakalei-Idate belongs and on whether Welaun is a Central Timor language or a dialect of Tetun. Hull groups Lakalei-Idate with Mambae, Kemak and Tokodede as Ramelaic; Grimes and Edwards (2026) place it in Timor-Babar, apart from Central Timor. This paper tests both against basic vocabulary from LexiRumah, under a protocol frozen before analysis and nine dated addenda. The pre-registered bootstrap rule leaves both questions unresolved. Bayesian analyses of 24 lects support Central Timor with Welaun (posterior 0.86, ln BF +10.3) and place Lakalei-Idate as its sister (0.91). No analysis gives lexical support to Hull's Fabronic, to Timor-Babar or to Eastern Timor, a result that survives consensus cognate coding, a retention classifier validated against experts, and word lists collected independently of Grimes and Edwards; Timor-Babar, defined by sound changes, is not refuted by this. On those independent lists Lakalei-Idate groups with Mambae, Kemak and Tokodede at 97\%. The arguments for its eastern placement do not discriminate once tested: the coded sound changes cost the same on either tree, and a borrowing simulation cannot separate the placements when borrowing runs both ways. Contact signals other than Welaun's Tetun loans do not survive calibration or multiple-testing correction, and the confirmatory predictions failed against Tetun. Lakalei-Idate's position remains formally unresolved but leans towards Central Timor, as an old member of a central Timor dialect chain. All data and code are public.\par
\vspace{2mm}
{\headingregular\itshape\color{ink2}Keywords:} Timor; Austronesian; subgrouping; computational phylogenetics; language contact; pre-registration\par}
\vspace{3mm}}}
\endgroup
\vspace{4mm}
